% 第十一部：2006 年後的發展與回應。出處：課程倉 W06_1019_導讀問答.md「本文發表後與本文論證相關的發展」一則；
% 原件在 1151/seminar/readings/W06/sources/ 與翻譯倉 JudicialReview/Quellen/raw/發展_20261007/。
\chapter{2006 年後的發展與回應}

\qaentry{第十一部之一}{作者的後續著作與論者的回應}

\begin{qaquestion}
本文發表後，作者本人與其他論者有哪些直接相關的著作？
\end{qaquestion}

本文第伍部分指出，法院內部的簡單多數決欠缺公平性的論證\pg{1391--1392}。作者 2014 年在 Yale Law Journal 第 123 卷發表〈Five to Four: Why Do Bare Majorities Rule on Courts?〉\nobreak，專以這個問題為題。\footcite{waldron2014}

Fallon 2008 年以〈The Core of an Uneasy Case for Judicial Review〉為題回應本文，\footcite{fallon2008} 第 4 週指定的 Fallon 專書第 1 章另在註 25 引本文 1387 至 1389 頁。\footcite[181 註 25]{fallon2018}

\qaentry{第十一部之二}{美國：本文所引兩案的後續}

\begin{qaquestion}
本文開篇的同性婚姻判決與多次引用的 Roe v.\ Wade，其後各有什麼發展？
\end{qaquestion}

本文以麻薩諸塞州 2003 年的 Goodridge 案開篇\pg{1348}。2015 年 6 月 26 日，聯邦最高法院在 Obergefell v.\ Hodges 以五比四判決，各州禁止同性婚姻或拒絕承認他州合法締結的同性婚姻，違反第十四修正案。\footcite{obergefell2015}

本文多次以 Roe v.\ Wade 為例：導論說支持選擇權的一方因擔心動搖 Roe，不願考慮反對違憲審查的論證\pg{1351}；第肆部分說 Roe 五十頁的意見書中處理墮胎道德問題的只有幾段\pg{1383}。2022 年 6 月 24 日，聯邦最高法院在 Dobbs v.\ Jackson Women's Health Organization 推翻 Roe 與 Casey，法院意見寫：

\begin{quote}
Roe and Casey must be overruled, and the authority to regulate abortion must be returned to the people and their elected representatives.\footcite{dobbs2022}
\end{quote}

\begin{readernote}
Dobbs 把墮胎的規範權交還各州的立法機關，與本文主張的決定場域相同；作成這項決定的仍是以多數決投票的大法官，可對照第伍部分對法院內部簡單多數決的批評\pg{1391--1392}。支持與反對 Roe 的兩方在 Dobbs 前後對違憲審查的態度，也可以拿 Charles Black 所說的 a sometime thing\pg{1351} 來對照。
\end{readernote}

\qaentry{第十一部之三}{紐西蘭與加拿大的變化}

\begin{qaquestion}
本文用來界定弱式與居間形式的紐西蘭、加拿大，其後制度與實務有何變化？
\end{qaquestion}

\subsection{紐西蘭}

本文寫作時，紐西蘭法院即使自行宣告法律與權利法案不相容，該宣告對立法程序也不生法律效果\pg{1356}。2018 年，紐西蘭最高法院在 Attorney-General v Taylor [2018] NZSC 104 確認高等法院有權作成不相容宣告，該案涉及 2010 年禁止受刑人投票的修正法。2022 年國會通過 New Zealand Bill of Rights (Declarations of Inconsistency) Amendment Act 2022，新增第 7A 條：上級法院（senior court，含高等法院、上訴法院與最高法院）就法律與權利法案不一致所為的宣告確定後，檢察總長應於眾議院第六個開議日以前提出通知，使眾議院知悉該宣告。\footnote{New Zealand Bill of Rights (Declarations of Inconsistency) Amendment Act 2022 (2022 No 45), \url{https://www.legislation.govt.nz/act/public/2022/45/en/latest}。}

\subsection{加拿大}

本文把加拿大計為強式審查，理由是儘管條款在實務上甚少援用\pg{1356--1357}。2019 年魁北克省為通過禁止部分公務員在職務上配戴宗教標誌的第 21 號法案，預先援用該條款；其後安大略省（2021、2022）、薩斯喀徹溫省（2023）等也相繼援用，涉及的事項包括第三方選舉廣告支出的限制（安大略 2021）、禁止加拿大公共雇員工會所屬教育工作者罷工（安大略 2022，旋即廢止）、未滿 16 歲學生在校更改名字與代名詞須得家長同意（薩斯喀徹溫 2023）等。\footnote{\url{https://www.lawnow.org/where-are-my-charter-rights/}；\url{https://www.cp24.com/news/canada/2025/09/17/alberta-defends-notwithstanding-clause-in-arguments-for-bill-21-supreme-court-case/}。}

\begin{readernote}
紐西蘭的不相容宣告自此對立法程序發生作用，在本文的分類上與英國的弱式審查距離縮短；加拿大歸類所依據的前提「甚少援用」，在 2019 年以後已有變化。
\end{readernote}
